%----------------------------------------------------------------------------------------------------
% START OF CHAPTER
%----------------------------------------------------------------------------------------------------
\chapter{Apr 27}
\paragraph{Topics}
\textit{Intro to the past tense. Sentences with the past tense. Grammar analysis.}

\section{Grammar} \label{0427-Grammar}
\subsection{Past tense}

In English, `I have spoken' and `I spoke' are grammatically equivalent in everyday use. In Italian,
everyday speech and writing uses the latter form, with \iti{essere} and \iti{avere} acting as helper verbs.
The past tense therefore uses a three-word structure: \textit{subject, helper-verb, past-participle}.

Broadly speaking, transitive verbs use \iti{avere} and intransitive verbs and verbs of motion use \iti{essere}.

\paragraph{Verbs that require \itb{essere}}
\begin{itemize}
    \item \iti{essere}
    \item \iti{stare}
    \item \iti{entrare}
    \item \iti{arrivare}
    \item \iti{venire}
    \item \iti{andare}
    \item \iti{uscire}
    \item \iti{partire}
\end{itemize}

To form the past participle, remove the ending of the verb (e.g., \iti{-are, -ere, -ire} and then add 
\iti{-ato, -uto, -ito} to match. If \iti{essere} is the helper verb, then change the ending to match gender and plurality. 

\begin{itemize}
    \item \iti{comprare} $ \rightarrow $ \iti{comprato}
    \item \iti{lavorare} $ \rightarrow $ \iti{lavorato}
    \item \iti{vedere}   $ \rightarrow $ \iti{visto}    
    \item \iti{arrivare} $ \rightarrow $ \iti{arrivato} or \iti{arrivata, arrivati} 
\end{itemize}

\begin{table}[ht]
\centering
    \begin{tabular}{ r | l }
         x   & Helper: \itb{avere} Verb: \itb{comprare}    \\
         \hline
        \itb{io}       & \iti{ho comprato}      \\
        \itb{tu}       & \iti{hai comprato}     \\
        \itb{lui/lei}  & \iti{ha comprato}      \\
        \itb{noi}      & \iti{abbiamo comprato} \\
        \itb{voi}      & \iti{avete comprato}   \\
        \itb{loro}     & \iti{hanno comprato}  
      \end{tabular}
    \caption{Past tense example: \iti{avere + comprare}
    \label{tab:tPastTenseAvereComprare}
\end{table}

\begin{table}[ht]
\centering
    \begin{tabular}{ r | l }
         x   & Helper: \itb{essere} Verb: \\itb{andate}    \\
         \hline
        io       & \iti{sono andato}   \\
        tu       & \iti{sei andato}    \\
        lui      & \iti{è andato}      \\
        lei      & \iti{è andata}      \\
        noi      & \iti{siamo andati}  \\
        voi      & \iti{siete andati}  \\
        loro     & \iti{sono andati}  
      \end{tabular}
    \caption{Past tense example: \iti{essere + andare}
    \label{tab:tPastTenseEssereAndare}
\end{table}

Unsurprisingly, some verbs are irregular and you just have to learn them.

\begin{table}[ht]
\centering
    \begin{tabular}{ r | l }
         Past participle   & Helper: \itb{avere} Irregular verbs   \\
         \hline
        \textbf{said, told}     & \iti{detto}   \\
        \textbf{made, did}      & \iti{fatto}   \\
        \textbf{read}           & \iti{letto}   \\
        \textbf{seen, saw}      & \iti{visto}   \\
        \textbf{put}            & \iti{messo}   \\
        \textbf{wrote, written} & \iti{scritto}  
      \end{tabular}
    \caption{Irregular past particples: all with \iti{avere}
    \label{tab:tIrregularPastParticiples}
\end{table}

\paragraph{Examples}
\begin{itemize}
  \item I'm happy because my wife made a really great dish of pasta last night. We watched a good move after dinner. \\
  \iti{Io sono felice perché mia moglie ha hatto un piatto di pasta molto buono ieri sera. Noi abbimamo vitso un buon film dopo cena.

  \item I said that I worked a lot today and am tired. I don't want to go to the restaurant for dinner. I prefer to stay at home. \\
  \iti{Ho detto che ho lavorato molto oggi e sono stanco. Io non voglio andare al ristorante per cena. Preferisco stare a casa
\end{itemize}

\section{Grammar analysis}
IO NON SONO COMPRANDO UNA CASA IN ITALIA PERCHE’ E’ MOLTO CARA
IO NON STO COMPRANDO UNA CASA IN ITALIA PERCHE’ E’ MOLTO CARA

IO NON POSSO PARLA ADESSO PERCHE’ IO STO LAVORANDO
IO NON POSSO PARLARE ADESSO PERCHE’ IO STO LAVORANDO

IO HO LA MACCHINA DI CLAUDIO (CORRETTO)

IO HO IL BICCHIERE DI VINO DI CLAUDIO (CORRETTO)

MIO LIBRO E’ MOLTO DIFFICILE
IL MIO LIBRO E’ MOLTO DIFFICILE

MIO FRATELLO E’ IN ITALIA (CORRETTO)

LA MIA AMICA E’ LA SORELLA DI CLAUDIO (CORRETTO)

IO MANGIO QUESTA PIATTO DI PASTA PERCHE E’ MOLTO BUONO
IO MANGIO QUESTA PIATTO DI PASTA PERCHE E’ MOLTO BUONO

IO VOGLIO ANDARE IN BANCA PERCHE’ HO BISOGNO DI CAMBIARE EURO (CORRETTO)

IO NON SONO FELICE PERCHE’ IO NON HO IL BIGLIETTO DEL TRENO (CORRETTO)
%----------------------------------------------------------------------------------------------------
% END OF CHAPTER
%----------------------------------------------------------------------------------------------------
